\section{Limitações do trabalho}

Apesar dos resultados obtidos, certas limitações devem ser consideradas no modelo de simulação desenvolvido neste trabalho.

Em primeiro lugar, o modelo orbital adotado baseia-se em hipóteses simplificadas, como a utilização das equações de \gls{hcw} para descrever o movimento relativo entre as espaçonaves. 
Esse modelo considera órbitas próximas e circularidade da órbita de referência, não incluindo efeitos perturbativos como achatamento da Terra (J2), arrasto atmosférico ou influência de terceiros corpos.

Outra limitação está relacionada à modelagem do sistema de controle. 
As leis de controle utilizadas são baseadas em controladores \gls{pd}, que, embora mostraram-se adequadas para análise inicial do comportamento do sistema, não representam necessariamente estratégias de controle mais avançadas que poderiam ser utilizadas em missões reais.

Além disso, a modelagem do \gls{mr} foi realizada a partir de parâmetros físicos obtidos a partir de um modelo CAD, sendo utilizada uma representação simplificada de sua estrutura dinâmica. 
Aspectos como flexibilidade estrutural dos elos e atrasos de atuadores não foram considerados no modelo.

Por fim, a simulação não inclui modelos detalhados de sensores e atuadores utilizados na determinação da posição relativa e da atitude das espaçonaves. 
A inclusão desses elementos poderia permitir uma representação mais próxima de um sistema embarcado real.
